% --- PREÁMBULO FINAL OPTIMIZADO PARA LA EDITORIAL UD ---
\usepackage[spanish,es-tabla,es-lcroman]{babel}
% Estilo editorial
\input{esConfigDir/esConfigFonts} % Definición de fuentes (lhseries, medseries, etc.)
\input{esConfigDir/esEstilo}

\usepackage{geometry}
\geometry{
	paperheight=24.0cm,
	paperwidth=17.0cm,
	textwidth=12.5cm,
	tmargin=2.5cm,
	bmargin=2.5cm,
	left=2.5cm
}


%%%%%%%%%%%%%%%%%%%%%%%%%%%%%%%%%%%%%%%%%%%%%%%%%%%%%%%%%%%%%%%%%%%%%%%%%%%%%%%%%%%%%%%%%%%%%%%%%
%%% BIBLIOGRAFÍA 
%%%%%%%%%%%%%%%%%%%%%%%%%%%%%%%%%%%%%%%%%%%%%%%%%%%%%%%%%%%%%%%%%%%%%%%%%%%%%%%%%%%%%%%%%%%%%%%%%
\usepackage[backend=biber,style=ieee,sorting=none,isbn=false,doi=true]{biblatex}
\addbibresource{references.bib}
\setlength{\bibitemsep}{1.5ex}

\DeclareQuoteStyle{spanish}{``}{''}{`}{'} % Reemplaza las comillas angulares por comillas dobles normales

\DefineBibliographyStrings{spanish}{
  url = {[En línea]. Disponible en},
  translator = {trad\adddot},
  bytranslator = {trad\adddot},
  version = {versión},
  and = {and},
}
% IEEE + español: delimitador final de lista de autores
\DeclareDelimFormat{finalnamedelim}{%
  \ifnumgreater{\value{liststop}}{2}{\finalandcomma\addspace}{}%
  \bibstring{and}\addspace}
\DefineBibliographyStrings{english}{%
  andothers = {et\addabbrvspace al\adddot}%
}

%\usepackage{showframe}% Similar al calderon, apagar al terminar de maquetar.
\usepackage[cam, width=215truemm, height=275truemm, lualatex, center]{crop} % Descomentar solo en producción final

% 5. AJUSTES FINALES DE FORMATO DEL PÁRRAFO
\usepackage[nopatch=footnote]{microtype} % Microtipografía avanzada
\usepackage{ragged2e} % Para control de alineación
\usepackage{authblk} % Manejo de autores (plantilla experimento)
\decimalpoint % Para usar la coma como separador decimal.
\setlength\parindent{0pt} % Sin sangría al inicio de párrafos, como pide la UD.
\newcommand{\normalfonts}{\fontsize{12pt}{14.4pt}\selectfont}
\let\normalsize\normalfonts

%%%%%%%%%%%%%%%%%%%%%%%%%%%%%%%%%%%%%%%%%%%%%%%%%%%%%%%%%%%%%%%%%%%%%%%%%%%%%%%%%%%%%%%%%%%%%%%%%
%%% HYPERREF (después de biblatex para que los enlaces de citas funcionen)
%%%%%%%%%%%%%%%%%%%%%%%%%%%%%%%%%%%%%%%%%%%%%%%%%%%%%%%%%%%%%%%%%%%%%%%%%%%%%%%%%%%%%%%%%%%%%%%%%
\usepackage[hidelinks,unicode,naturalnames,bookmarksopen,bookmarksdepth=2,breaklinks=true,hyperfootnotes=false]{hyperref}
\usepackage{bookmark}
\urlstyle{same}
\newcommand{\secref}[1]{``\nameref{#1}''}
\renewcommand{\textcite}{\cite}

% ÍNDICES Y SUBARCHIVOS
%\usepackage{makeidx}
%\usepackage{subfiles}

%%%%%%%%%%%%%%%%%%%%%%%%%%%%%%%%%%%%%%%%%%%%%%%%%%%%%%%%%%%%%%%%%%%%%%%%%%%%%%%%%%%%%%%%%%%%%%%%%
%%% CORNISAS Y ESTILO DE PÁGINA
%%%%%%%%%%%%%%%%%%%%%%%%%%%%%%%%%%%%%%%%%%%%%%%%%%%%%%%%%%%%%%%%%%%%%%%%%%%%%%%%%%%%%%%%%%%%%%%%%

\usepackage{emptypage}

% Esto asegura que la página izquierda de un capítulo nuevo esté vacía (sin cornisas)
\renewcommand{\cleardoublepage}{%
  \clearpage\ifodd\value{page}\else\hbox{}\thispagestyle{empty}\newpage\fi%
}

%%%%%%%%%%%%%%%%%%%%%%%%%%%%%%%%%%%%%%%%%%%%%%%%%%%%%%%%%%%%%%%%%%%%%%%%%%%%%%%%%%%%%%%%%%%%%%%%%
%%% PAQUETES FUNCIONALES Y DE CONTENIDO
%%%%%%%%%%%%%%%%%%%%%%%%%%%%%%%%%%%%%%%%%%%%%%%%%%%%%%%%%%%%%%%%%%%%%%%%%%%%%%%%%%%%%%%%%%%%%%%%%

% PAQUETES DE MATEMÁTICAS Y FÍSICA
\usepackage{amsmath}
\usepackage{amssymb}
\usepackage{amsfonts}
\usepackage{amsthm}
\usepackage{bm}
\usepackage{siunitx} % \num, \SI — antes de physics para que siunitx no avise al cargar
\usepackage{physics}
\AtBeginDocument{\RenewCommandCopy\qty\SI} % physics redefine \qty; restaurar definición de siunitx
% El aviso de siunitx+physics usa expl3 (no \WarningFilter); silenciar solo ese mensaje:
% https://tex.stackexchange.com/q/681700
\ExplSyntaxOn
\msg_redirect_name:nnn { siunitx } { physics-pkg } { none }
\ExplSyntaxOff
\usepackage{fontawesome}
\usepackage{fancyvrb}
\usepackage{fp}
\usepackage{calculus}
\usepackage{calculator}
\usepackage{mathtools}
\usepackage{bigints}
\usepackage{resizegather}

% PAQUETES DE GRÁFICOS E IMÁGENES
\usepackage{graphicx}
\usepackage[export]{adjustbox}
\graphicspath{{./Fig}{./Img}{./Imgn}}

% COLOR Y LAYOUT
\usepackage{xcolor}
\definecolor{orcidlogocol}{HTML}{A6CE39}
\usepackage{academicons}
\usepackage{multicol}
\usepackage{esvect}
\usepackage{relsize}
\usepackage{pifont}
\usepackage{etoolbox}
\usepackage{tcolorbox}
\tcbuselibrary{breakable,skins}

% Default options for tcolorbox (Instrucción Mathematica and similar)
\tcbset{
  enhanced,
  before skip=1em,
  after skip=1em,
}

\newtcolorbox{graybox}[1][]{
  enhanced,
  before skip=1em,
  after skip=1em,
  colback=gray!5,
  colframe=gray!80,
  title={\bf #1}
}

\usepackage{float}
\usepackage{caption}
\usepackage{array}
\usepackage[absolute,overlay]{textpos}
\usepackage{listings}
\usepackage{afterpage}
\usepackage{placeins}
\usepackage{wrapfig}
\usepackage{subcaption}
\usepackage{needspace}
\usepackage{pdfpages}
\usepackage{answers}
\usepackage{./Estilos/listofanswers}
\usepackage{empheq}

%%%%%%%%%%%%%%%%%%%%%%%%%%%%%%%%%%%%%%%%%%%%%%%%%%%%%%%%%%%%%%%%%%%%%%%%%%%%%%%%%%%%%%%%%%%%%%%%%
%%% CONFIGURACIONES Y COMANDOS
%%%%%%%%%%%%%%%%%%%%%%%%%%%%%%%%%%%%%%%%%%%%%%%%%%%%%%%%%%%%%%%%%%%%%%%%%%%%%%%%%%%%%%%%%%%%%%%%%

% COMANDOS PERSONALIZADOS Y ENTORNOS
\newcommand{\letraverbatim}{\fontsize{9pt}{11.5pt}\selectfont}
\newcommand{\ds}{\displaystyle}
\newcommand{\bi}{\int}
\newcommand{\la}{\langle}
\newcommand{\ra}{\rangle}
\newcommand{\ei}[1]{\mathbf{e}_{#1}}
\newcommand{\ve}[1]{\mathbf{#1}}
\DeclareMathOperator{\arcsinh}{arcsinh}
\DeclareMathOperator{\arccosh}{arccosh}
\DeclareMathOperator{\arctanh}{arctanh}
\DeclareMathOperator{\arccoth}{arccoth}
\DeclareMathOperator{\arcsech}{arcsech}
\DeclareMathOperator{\arccsch}{arccsch}
\DeclareMathOperator{\rot}{rot}
\newcommand{\vi}{\mathbf{i}}
\newcommand{\vj}{\mathbf{j}}
\newcommand{\vk}{\mathbf{k}}
\newcommand{\brake}{\\[4pt]}
% Cajas tcolorbox y atajos (preambulo_experimento)
\newcommand{\ptc}{\tcbset{colback=gray!10!white,colframe=black!85!blue}\begin{tcolorbox}[breakable,enhanced,title=\bf Parametrizando el sólido]}
\newcommand{\ptr}{\tcbset{colback=gray!10!white,colframe=black!85!blue}\begin{tcolorbox}[breakable,enhanced,title=\bf Parametrizando la región de Integración]}
\newcommand{\tc}{\begin{tcolorbox}[breakable,enhanced,title=\bf Instrucción en {\tt Mathematica}]}
\newcommand{\etc}{\end{tcolorbox}}
% PSTricks: usar solo con flujo compatible (p. ej. auto-pst-pdf); no cargado aquí por LuaLaTeX.

% Símbolo de final de ejemplo, demostración, etc.
% Referencia: https://www.ams.org/arc/tex/howto/qed/extra-qed.pdf
\newcommand{\xqed}[1]{
  \leavevmode\unskip\penalty9999 \hbox{}\nobreak\hfill
\quad\hbox{\ensuremath{#1}}}
\newcommand{\finalm}{\xqed{\square}}
\newcommand{\finalmf}{\vspace*{-\baselineskip}\xqed{\square}}
\newcommand{\finaleq}[1][-2.25\baselineskip]{\relax\par\vspace*{#1}\hfill$\square$\par\vspace*{\belowdisplayskip}\relax}
\let\final\finalm

% ENTORNOS DE TEOREMA Y REDEFINICIONES
\newtheorem{teorema}{Teorema}[section]
\newtheorem{prop}{Proposición}[section]
\newtheorem{defi}{Definición}[section]
\newtheorem{ejer}{Ejercicio}[section]
\newtheorem{ejemplo}{Ejemplo}[section]
\let\sen\sin
\renewcommand{\proofname}{Demostración}
% Demostración (macro propia)
\newcommand{\iniciodemostracion}{\noindent\textbf{Demostración.}\hspace{0.5em}}
\newcommand{\findemostracion}{\final}

\renewenvironment{proof}
{\vspace{0.5em}\iniciodemostracion\normalfont\ignorespaces}
{\findemostracion\vspace{0.5em}\par} %

\renewcommand*{\tablename}{\bf Tabla} % Estilo UD
\renewcommand*{\figurename}{\bf Figura} % Estilo UD

% AJUSTES DE ÍNDICE Y PROFUNDIDAD DE SECCIONES
\counterwithout{footnote}{chapter}
\setcounter{tocdepth}{3}
\setcounter{secnumdepth}{3}
\renewcommand*\contentsname{\lhseries Contenido}
\renewcommand{\cfttoctitlefont}{\hfill\Huge}
\renewcommand{\cftdot}{}

\makeindex

\raggedbottom
